% --- 1. INFORMAZIONI GENERALI ---
\section{Informazioni Generali}
\begin{itemize}
	\item \textbf{Luogo/Piattaforma:} Server Discord\textsubscript{G}
	\item \textbf{Data:} 2026-04-03
	\item \textbf{Orario di inizio:} 14:30
	\item \textbf{Orario di fine:} 16:00
	\item \textbf{Responsabile\textsubscript{G}:} Sofia De Blasi
	\item \textbf{Scriba:} Sofia De Blasi
\end{itemize}

\vspace{0.5cm}
\textbf{Partecipanti presenti:}
\begin{table}[ht]
	\centering
	\renewcommand{\arraystretch}{1.2}
	\begin{tabularx}{\textwidth}{X l}
		\toprule
		\textbf{Nome e Cognome} & \textbf{Matricola} \\
		\midrule
		Sofia De Blasi & 2111014 \\
		Roberto Mariano Doroftei & 2111031 \\
		Filippo Idiometri & 2035617 \\
		Francesco Marcon & 2110990 \\
		Armando Moda Scarati & 2082864 \\
		Anton Ricardo Rupi & 2054304 \\
		\bottomrule
	\end{tabularx}
\end{table}

\vspace{0.5cm}
\textbf{Partecipanti assenti:}
\begin{table}[ht]
	\centering
	\renewcommand{\arraystretch}{1.2}
	\begin{tabularx}{\textwidth}{X l}
		\toprule
		\textbf{Nome e Cognome} & \textbf{Matricola} \\
		\midrule
		Nessuno & - \\
		\bottomrule
	\end{tabularx}
\end{table}

% --- 2. ORDINE DEL GIORNO ---
\section{Ordine del Giorno}
Gli argomenti previsti per la riunione odierna sono:
\begin{enumerate}
	\item Pianificazione della documentazione da produrre nella fase RTB\textsubscript{G}:
	      Analisi dei Requisiti\textsubscript{G}, Piano di Progetto\textsubscript{G}, Piano di Qualifica\textsubscript{G},
	      Norme di Progetto\textsubscript{G}.
	\item Scelta delle tecnologie per il progetto.
	\item Aggiornamento delle convenzioni operative: gestione di task\textsubscript{G},
	      branch\textsubscript{G}, commit\textsubscript{G} e pull request\textsubscript{G}.
	\item Revisione\textsubscript{G} e aggiornamento del template\textsubscript{G} documentale.
\end{enumerate}

\newpage
% --- 3. RESOCONTO ---
\section{Resoconto della Riunione}

\subsection{Pianificazione della Documentazione RTB}

Il gruppo ha identificato i principali documenti da produrre nella fase RTB
e, dopo una panoramica generale del loro scopo e contenuto, ha definito le
modalità di avvio per ciascuno.

\paragraph{Analisi dei Requisiti} 
Il documento costituisce la base della fase RTB e include la definizione dei casi d'uso\textsubscript{G}, con relativa
rappresentazione tramite diagrammi UML\textsubscript{G}. Il gruppo ha stabilito di studiare
preventivamente il materiale del corso prima di avviare la stesura.

\paragraph{Piano di Progetto}
Il documento regolerà la gestione del progetto tramite sprint\textsubscript{G}, includendo
le attività di sprint planning\textsubscript{G}, sprint review\textsubscript{G} e sprint retrospective\textsubscript{G},
la rotazione dei ruoli e il tracciamento delle attività. Il gruppo ha
stabilito di avviarne la stesura parallelamente all'Analisi dei Requisiti.

\paragraph{Piano di Qualifica}
Il documento definirà le metriche quantitative per la qualità di processo
e di prodotto. Il gruppo ha stabilito di rimandarne la stesura a una
comprensione più approfondita delle lezioni di riferimento.

\paragraph{Norme di Progetto}
Il documento continuerà a essere aggiornato in modo incrementale secondo
il principio just-in-time\textsubscript{G} già adottato. In questa sessione sono state
definite le nuove convenzioni operative descritte nella
Sezione~\ref{sec:convenzioni}.

\subsection{Scelta delle Tecnologie}
\label{sec:tecnologie}

Il gruppo ha preso atto delle tecnologie obbligatorie indicate dal capitolato
(HTML, CSS, Markdown e API\textsubscript{G} compatibili con le specifiche OpenAI) e ha
stabilito di rimandare la scelta delle tecnologie complementari a una fase
successiva, da svolgere contestualmente all'avanzamento dell'Analisi dei
Requisiti. La definizione della technology baseline\textsubscript{G}\textsubscript{G} — obiettivo della
revisione RTB — sarà avviata dopo il completamento del 20--30\% dell'Analisi
dei Requisiti.
 
\subsection{Aggiornamento delle Convenzioni Operative}\label{sec:convenzioni}
 
Il gruppo ha ridefinito il flusso di lavoro per la gestione di task, branch,
commit e pull request, introducendo le seguenti convenzioni aggiornate.
 
\paragraph{Versionamento Documentale.}
Tutti i documenti partono dalla versione \texttt{0.0.0}.
 
\paragraph{Nomenclatura dei Task e dei Branch.}
Il nome del task, e il nome del branch corrispondente, seguono la convenzione:
 
\begin{center}
	\texttt{<id>-<[gestionale|tecnico]>\_<[esterno|interno]>\_<nomeDoc>}
\end{center}
 
Il branch viene creato direttamente dalla pagina della issue\textsubscript{G} su GitHub
(sezione \textit{Development}) e recuperato in locale con:
 
\begin{center}
	\texttt{git checkout -b <nome-branch> origin/<nome-branch>}
\end{center}
 
\paragraph{Collegamento Commit -- Issue.}
Ogni commit relativo a una issue deve includere il riferimento all'identificativo
della issue nel messaggio:
 
\begin{center}
	\texttt{git commit -m "\#<id> <messaggio>"}
\end{center}
 
\paragraph{Flusso di Lavoro Aggiornato.}
Il flusso operativo da seguire per ogni task è il seguente:
 
\begin{enumerate}
	\item Il Responsabile crea il task tramite Google Form\textsubscript{G} durante la riunione.
	\item Il Redattore\textsubscript{G} del verbale di riunione inserisce il riferimento alla issue nel verbale
	      (tabella delle azioni).
	\item L'Assegnatario di una task crea il branch direttamente dalla issue su GitHub e
	      lo recupera in locale.
	\item Ad ogni commit significativo, l'Assegnatario include il riferimento
	      alla issue nel messaggio di commit.
	\item Al termine del lavoro, l'Assegnatario apre la pull request e collega
	      la issue corrispondente.
	\item Il Verificatore\textsubscript{G} effettua la revisione; l'Assegnatario (non il Verificatore)
	      esegue il merge\textsubscript{G} dopo l'approvazione.
	\item Il branch viene eliminato dopo il merge.
\end{enumerate}
 
\paragraph{Automazione Google Form.}
È stato discusso di configurare il Google Form per la creazione dei task con
campi precompilati che rispettino automaticamente la convenzione di nomenclatura
(es.\ tipologia gestionale/tecnico, categoria esterno/interno, nome documento).
Questo consentirà di ridurre gli errori manuali nella creazione dei task.
 
\subsection{Revisione del Template Documentale}
 
Il gruppo ha identificato le seguenti necessità di aggiornamento del template
condiviso:
 
\begin{itemize}
	\item Aggiornamento della configurazione dei link (\verb|\href{url}{\underline{text}}|)
	      per garantire coerenza visiva tra i documenti.
	\item Revisione della struttura delle tabelle (divisione con "Continua nella pagina successiva", font registro modifiche,
	      allineamento colonne) per uniformità e corretto comportamento in
	      caso di contenuto lungo.
	\item Aggiornamento della tabella delle azioni (sezione To-Do dei verbali)
	      per includere il codice identificativo della issue GitHub come link
	      cliccabile, con la sintassi: 		  
	      \verb|\href{url}{\underline{\#id}}|.
	\item Configurazione dell'account \textit{Synergo Unit} come token per
	      le chiamate API delle automazioni del Google Form.
\end{itemize}

\newpage

% --- 4. DECISIONI PRESE ---
\section{Decisioni Prese}

\renewcommand{\arraystretch}{1.3}
\vspace{-0.3cm}
\noindent
\begin{longtable}{c p{12cm}}
	\toprule
	\textbf{Codice} & \textbf{Decisione} \\
	\midrule
	\endfirsthead

    \toprule
    \textbf{Codice} & \textbf{Decisione} \\
    \midrule
    \endhead

    \midrule
    \multicolumn{2}{r}{\textit{Continua nella pagina successiva}} \\
    \midrule
    \endfoot

	\bottomrule
    \endlastfoot

	\textbf{VI-7.1} & Tutti i documenti partono dalla versione \texttt{0.0.0}. \\[4pt]
	\textbf{VI-7.2} & I branch assumono il nome dal task corrispondente secondo
					la convenzione \texttt{<id>-<gestionale|tecnico>\_<esterno|interno>\_<nomeDoc>}
					e vengono creati direttamente dalla issue su GitHub. \\[4pt]
	\textbf{VI-7.3} & Ogni commit relativo a una issue include il riferimento
					\texttt{\#id} nel messaggio. \\[4pt]
	\textbf{VI-7.4} & La pull request viene aperta dal Redattore con collegamento
					alla issue; il merge è eseguito dal Redattore dopo
					l'approvazione del Verificatore. \\[4pt]
	\textbf{VI-7.5} & Il Google Form per la creazione dei task sarà configurato
					con campi precompilati per la nomenclatura dei branch. \\[4pt]
	\textbf{VI-7.6} & La technology baseline sarà avviata dopo il completamento
					del 20--30\% dell'Analisi dei Requisiti. \\[4pt]
	\textbf{VI-7.7} & Il template documentale sarà aggiornato per coerenza
					di tabelle, link e versionamento. \\
\end{longtable}
 
% --- 5. AZIONI DA SVOLGERE ---
\section{Azioni da Svolgere (To-Do)}

\renewcommand{\arraystretch}{1.3}
\begin{longtable}{c c p{4.5cm} p{3cm} l}
    \toprule
    \textbf{Codice} & \textbf{Rif.} & \textbf{Descrizione Task} & \textbf{Assegnatario} & \textbf{Scadenza} \\
    \midrule
    \endfirsthead

    \toprule
    \textbf{Codice} & \textbf{Rif.} & \textbf{Descrizione Task} & \textbf{Assegnatario} & \textbf{Scadenza} \\
    \midrule
    \endhead

    \midrule
    \multicolumn{5}{r}{\textit{Continua nella pagina successiva}} \\
    \midrule
    \endfoot

    \bottomrule
    \endlastfoot

    \href{https://github.com/Synergo-Unit-UniPD/Documentary-Repo/issues/93}{\underline{\#93}} & VI-9.7
        & {Aggiornamento configurazione template: tabelle (divisione,
           versionamento, allineamento), href e coerenza visiva.
           (Revisore: Filippo)}
        & Armando Moda Scarati & 2026-04-07 \\[4pt]

    \href{https://github.com/Synergo-Unit-UniPD/Documentary-Repo/issues/94}{\underline{\#94}} & VI-9.1
        & {Stesura verbale interno della riunione odierna.
           (Revisore: Ricardo)}
        & Sofia De Blasi & 2026-04-07 \\[4pt]

    \href{https://github.com/Synergo-Unit-UniPD/Documentary-Repo/issues/95}{\underline{\#95}} & VI-9.5
        & {Configurazione campi precompilati Google Form per nomenclatura
           task e branch.}
        & Francesco Marcon & 2026-04-07 \\[4pt]

    \href{https://github.com/Synergo-Unit-UniPD/Documentary-Repo/issues/96}{\underline{\#96}} & --
        & {Studio struttura e contenuto dell'Analisi dei Requisiti
           (casi d'uso, pattern Cardin, UML).}
        & Filippo Idiometri, Roberto M.\ Doroftei, Armando M.\ Scarati & 2026-04-07 \\[4pt]

    \href{https://github.com/Synergo-Unit-UniPD/Documentary-Repo/issues/97}{\underline{\#97}} & --
        & {Studio struttura e contenuto del Piano di Progetto.}
        & Anton Ricardo Rupi, Francesco Marcon, Sofia De Blasi & 2026-04-07 \\[4pt]
	
	\href{https://github.com/Synergo-Unit-UniPD/Documentary-Repo/issues/99}{\underline{\#99}} & --
	& {Aggiornamento sito web con nuova sezione RTB.
		(Revisore: Ricardo)}
	& Sofia De Blasi & 2026-04-07 \\
\end{longtable}

 
% --- 6. PROSSIMA RIUNIONE ---
\section{Prossima Riunione}
\begin{itemize}
	\item \textbf{Data:} 2026-04-07
	\item \textbf{Ora:} 14:30
	\item \textbf{OdG preliminare\textsubscript{G}:}
	\begin{itemize}
		\item Ragionare sulla gestione di task non strettamente legate
		      alla repository\textsubscript{G} e con più assegnatari (es.\ studio individuale (tempo di palestra),
		      attività di formazione, task che non producono cambiamenti nella repo).
		\item Discussione sulla possibilità di automatizzare l'apertura
		      dei branch con checkbox sì/no.
		\item Aggiornamento sull'avanzamento delle attività assegnate.
	\end{itemize}
\end{itemize}